% Copyright IBM Corp. All Rights Reserved.
% SPDX-License-Identifier: Apache-2.0
%
% Titan PCS -- the cryptography as implemented in
% token/core/zkatdlog/nogh/v1/crypto/titan.
%
% Self-contained: standard classes only, no dependency on the Titan paper's
% private macro set. Build with pdflatex.

\documentclass[11pt,a4paper]{article}

\usepackage[margin=1in]{geometry}
\usepackage{amsmath,amssymb,amsthm}
\usepackage{booktabs}
\usepackage{enumitem}
\usepackage[hidelinks]{hyperref}

\newcommand{\F}{\mathbb{F}}
\newcommand{\G}{\mathbb{G}}
\newcommand{\dom}{\mathcal{L}}
\newcommand{\bin}[1]{\langle #1 \rangle}
\newcommand{\eq}{\widetilde{\mathrm{eq}}}
\newcommand{\Com}{\mathsf{Com}}
\newcommand{\paper}[1]{\textsf{[Titan, #1]}}

\theoremstyle{definition}
\newtheorem{remark}{Remark}

\title{Titan PCS: the cryptography as implemented}
\author{Panurus \textemdash\ \texttt{crypto/titan}}
\date{}

\begin{document}
\maketitle

\begin{abstract}
\noindent
This note describes the Titan polynomial commitment scheme \emph{as implemented} in
the Panurus \texttt{crypto/titan} package: what is committed, what the evaluation
proof consists of, and which check makes the opening binding. It follows the Titan
paper's construction and cites its sections, but it documents the code, so it also
records where the implementation deliberately departs from the paper and what it does
not yet do. It is not a security proof; for that, see the paper.
\end{abstract}

\section{Setting and notation}

Work over BLS12-381: $\F$ is the scalar field $\F_r$, and $\G$ is $\mathbb{G}_1$
written additively. A multilinear polynomial in $m$ variables is given by its
$2^m$ evaluations on $\{0,1\}^m$; we write $\widetilde{f}$ for the multilinear
extension and $\bin{i}$ for the binary decomposition of $i$ into the relevant number
of bits. As usual,
\[
  \eq(\vec X, \vec Y) = \prod_{j} \bigl( X_j Y_j + (1-X_j)(1-Y_j) \bigr),
  \qquad
  \widetilde{f}(\vec\alpha) = \sum_{\vec x \in \{0,1\}^m} \eq(\vec\alpha, \vec x)\, \widetilde{f}(\vec x).
\]
A \emph{group multilinear} is the same object with coefficients in $\G$; its
evaluation is a group element, computed by multi-scalar multiplication.

$\dom \subset \F^\ast$ is a smooth multiplicative subgroup of order $2^d$, used as the
Reed--Solomon evaluation domain. The rate is $\rho = 2^{-\mathsf{LogRate}}$, and
$\mathsf{LogRate} = 3$ by default, so $\rho = 1/8$.

\begin{remark}[a notation collision worth flagging]
The paper writes the committed matrix as $p \times q$ with variable blocks
$\ell_p, \ell_q$, and separately uses $m_1 := |\mathcal{K}_1|$ for a \emph{Merkle leaf
count} \paper{\S\,titan-commit}. The implementation's \texttt{Split.M1} is
\emph{unrelated} to that $m_1$: it is the number of \emph{column} variables of the
matrix, the paper's $\ell_q$. The leaf count appears in the code as
\texttt{Commitment.NumLeaves}.
\end{remark}

\section{The two-tier commitment}
\label{sec:commit}

Following \paper{\S\,titan-commit}, a scalar multilinear $\widetilde{f}$ in $m$
variables is read as a matrix. Let $\mathsf{M1}$ be the number of column variables and
$\mathsf{RowVars} = m - \mathsf{M1}$ the number of row variables, giving a
$2^{\mathsf{RowVars}} \times 2^{\mathsf{M1}}$ matrix and the decomposition
\begin{equation}
  \widetilde{f}(\vec X, \vec Y) \;=\; \sum_{i=0}^{2^{\mathsf{RowVars}}-1}
    \eq\bigl(\vec X, \bin{i}\bigr) \cdot \widetilde{f}\bigl(\bin{i}, \vec Y\bigr),
  \label{eq:decompose}
\end{equation}
where $\vec X$ ranges over the row variables and $\vec Y$ over the column variables.

\paragraph{Tier 1: Pedersen rows.}
With generators $g_0, \dots, g_{2^{\mathsf{M1}}-1} \in \G$, each row is committed as
\begin{equation}
  G_i \;=\; \Com\bigl(\widetilde{f}(\bin{i}, \cdot)\bigr)
        \;=\; \sum_{j=0}^{2^{\mathsf{M1}}-1} \widetilde{f}\bigl(\bin{i}, \bin{j}\bigr) \cdot g_j .
  \label{eq:rowcommit}
\end{equation}
This is one MSM of length $2^{\mathsf{M1}}$ per row. The $2^{\mathsf{RowVars}}$ results
are the coefficients of a \emph{group} multilinear
\begin{equation}
  \widetilde{G}(\vec X) \;=\; \sum_i \eq\bigl(\vec X, \bin{i}\bigr) \cdot G_i
  \label{eq:Gmle}
\end{equation}
in $\mathsf{RowVars}$ variables. Everything downstream is about $\widetilde G$; the
scalar polynomial appears again only in leg~2 of the evaluation proof.

\paragraph{Tier 2: encode and Merkle-commit.}
$\widetilde G$ is Reed--Solomon encoded over $\dom$ of size
$2^{\mathsf{RowVars} + \mathsf{LogRate}}$ along the \emph{power curve}: the codeword at
$x \in \dom$ is
\begin{equation}
  \widehat{G}(x) \;=\; \widetilde{G}\bigl(x, x^2, x^4, \dots, x^{2^{\mathsf{RowVars}-1}}\bigr),
  \label{eq:powercurve}
\end{equation}
computed by a butterfly recursion \paper{\S\,whir-oracle-algo}. This codeword is
\emph{not} committed: see the remark at the end of this section.

\paragraph{The coset oracle.}
The folding phase queries a \emph{second}, differently shaped commitment. For
$\ell$ folding rounds, a coset indexed by $y$ is
\begin{equation}
  \mathcal{C}_y \;=\; \bigl\{\, \widetilde{G}(\vec b,\, y, y^2, y^4, \dots) \;:\; \vec b \in \{0,1\}^{\ell} \,\bigr\},
  \label{eq:coset}
\end{equation}
boolean in the first $\ell$ coordinates and on the power curve in the rest. It is built
per slice: slice $\vec b$ is encoded over the folded domain $\dom^{2^\ell}$, and leaf
$y$ gathers index $y$ across the $2^\ell$ resulting codewords. These leaves form a
second Merkle tree, \texttt{Commitment.Cosets}.

These leaves form the Merkle tree, \texttt{Commitment.Cosets}, and it is the
\emph{only} commitment in the scheme.

\begin{remark}[one root, not two]
\label{rem:oneroot}
\eqref{eq:powercurve} and \eqref{eq:coset} are genuinely different objects --- a point
of the flat codeword versus a set of slice evaluations --- and an earlier version of
this implementation Merkle-committed \emph{both}, on the reasoning that a verifier
confusing them should fail loudly. But nothing ever verified against the flat root:
the folding phase of \S\ref{sec:fold} is the only stage that opens the oracle, and it
opens cosets. The flat encoding is a computational stepping stone to \eqref{eq:coset},
not a commitment anyone opens. Its tree is therefore no longer built, and
\texttt{Commitment.Root} is retained as a reserved, always-nil field --- an exported
root that binds nothing is a hazard, not a safeguard.
\end{remark}

\section{The evaluation proof}
\label{sec:eval}

To open $\widetilde f(\vec\alpha) = \sigma$, split $\vec\alpha$ at $\mathsf{M1}$ into
$\vec\alpha_{\mathrm{col}}$ (the first $\mathsf{M1}$ coordinates) and
$\vec\alpha_{\mathrm{row}}$. Substituting into \eqref{eq:decompose},
\[
  \widetilde f(\vec\alpha)
  = \sum_i \eq(\vec\alpha_{\mathrm{row}}, \bin i) \cdot \widetilde f(\bin i, \vec\alpha_{\mathrm{col}}),
\]
which the proof establishes in two legs \paper{\S\,titan-eval-proof}.

\paragraph{Leg 1 (group): $\widetilde G(\vec\alpha_{\mathrm{row}})$.}
A group sum-check over $\mathsf{RowVars}$ variables proves
$\sum_{\vec x} \eq(\vec\alpha_{\mathrm{row}}, \vec x)\, \widetilde G(\vec x) = \Sigma$
for a claimed $\Sigma \in \G$ \paper{\S\,group-sum-check}. Round messages are group
elements; the prover splits the work between an MSM part and a folklore part at a
tunable point.

\paragraph{Leg 2 (scalar): the partial evaluation.}
Leg~1 reduces to a claim about $\widetilde G$ at a transcript-derived point, which by
\eqref{eq:rowcommit} is a Pedersen commitment to a row combination. Leg~2 proves that
this commitment opens to the scalar vector
$\widetilde f(\cdot, \vec\alpha_{\mathrm{col}})$ via the CSP inner-product argument
\paper{\S\,partial-eval-csp}, and the linear form it checks pins $\sigma$
\paper{\S\,check-C}. This leg is \emph{not} folded here (see
\S\ref{sec:notdone}), so its cost is linear in $2^{\mathsf{M1}}$ --- which is why the
implementation keeps the column half small.

\begin{remark}[$\eq$ factorizes over \emph{any} split, and that is a trap]
\label{rem:eqtrap}
For any cut point, $\eq(\vec\alpha,\vec x)$ factorizes into a column part and a row
part. So a prover and a verifier that both split $\vec\alpha$ at the \emph{wrong} place
--- or at the same wrong place --- produce a proof that verifies against itself, for a
different polynomial. Acceptance cannot detect this. The implementation therefore
records the split in the commitment (\texttt{ColVars}) and cross-checks it, and its
tests assert $\sigma$ against an independent reference evaluator rather than relying on
a round trip.
\end{remark}

\section{The folding phase, and what makes the opening binding}
\label{sec:fold}

Legs 1 and 2 reduce the claim to a statement about a polynomial \emph{the verifier has
never queried}. Without more, a prover could commit to one polynomial and run the
reduction over another. The folding phase closes that gap.

Over $\ell$ rounds the claim is folded, one variable per round, against transcript
challenges. After the last round the prover sends the \emph{reduced} group multilinear
--- all $2^{\mathsf{RowVars}-\ell}$ coefficients, in plain --- together with $Q$
consistency queries, each a coset \eqref{eq:coset} with its Merkle path.

The verifier performs three checks:
\begin{enumerate}[label=(\arabic*),leftmargin=*]
  \item \textbf{Round consistency.} Each round message sums to the claim the previous
        round left.
  \item \textbf{Reduced claim.} $\langle \mathsf{Reduced},\, \eq(\vec r)\rangle$ equals
        the residual claim the folding rounds left.
  \item \textbf{Consistency queries.} For each of the $Q$ sampled indices: the coset
        opens under \texttt{Commitment.Cosets} (the only root; Remark~\ref{rem:oneroot}),
        \emph{and} it folds to the reduced
        polynomial's codeword at that index --- a single $\eq$ dot product,
        $\langle \mathcal{C}_y, \eq(\vec r)\rangle$.
\end{enumerate}

Checks (1) and (2) make the proof internally consistent. \textbf{Check (3) is the one
that makes it binding}: it is the only step that touches the committed oracle, and it
is what rejects a prover who folded a polynomial other than the one committed. A
convenient consequence of choosing the coset dimension equal to $\ell$ is that one
primitive --- an $\eq$ dot product --- serves both (2) and (3).

\paragraph{Query count.}
The query term dominates the soundness error; the folding rounds' Schwartz--Zippel term
is $\approx 2^{-252}$ and the $|\dom|/|\F|$ term $\approx 2^{-240}$. Hence
\begin{equation}
  Q \;=\; \left\lceil \frac{\lambda}{\log_2(1/\rho)} \right\rceil
  \;=\; \left\lceil \frac{128}{3} \right\rceil \;=\; 43
  \qquad (\lambda = 128,\ \rho = 1/8),
  \label{eq:queries}
\end{equation}
under the \emph{capacity} bound \paper{Thm.\,titan-pcs}. Capacity is
\textbf{conjectured}; the provable bound is Johnson's radius, which charges
$\tfrac{1}{2}\log_2(1/\rho)$ bits per query and so needs $86$ queries for the same
$\lambda$. The implementation defaults to capacity, matching the paper's cost analysis
\paper{\S\,titan-pcs-costs}, and exposes Johnson as a selectable regime.

\paragraph{Number of folding rounds.}
The proof carries $Q \cdot 2^{\ell}$ group elements of cosets plus
$2^{\mathsf{RowVars}-\ell}$ coefficients of the reduced polynomial, so the size-optimal
$\ell$ minimizes $Q\,2^{\ell} + 2^{\mathsf{RowVars}-\ell}$. Because the coset term
carries the factor $Q$, this optimum lies \emph{below} the paper's
$\mathsf{RowVars}/2 - 1$: at $\mathsf{RowVars} = 12$ it is $3$, not $5$. The
implementation computes it rather than fixing it, and leaves $\ell$ settable.

\section{Parameters and admissible sizes}

\begin{center}
\begin{tabular}{@{}lll@{}}
\toprule
parameter & default & meaning \\
\midrule
$\lambda$              & $128$          & target soundness in bits \\
$\mathsf{LogRate}$     & $3$            & $\rho = 1/8$ \\
regime                 & capacity       & conjectured; Johnson is provable \\
$Q$                    & $43$           & consistency queries, from \eqref{eq:queries} \\
$\ell$                 & size-optimal   & folding rounds \\
$\mathsf{M1}$          & $\lfloor m/2 \rfloor$ & column variables of the matrix \\
\bottomrule
\end{tabular}
\end{center}

Two structural constraints bound which $m$ are usable.

\emph{Even row half.} The coset layout halves the row variables exactly, so
$\mathsf{RowVars} = m - \mathsf{M1}$ must be even. Under the balanced split
$\mathsf{M1} = \lfloor m/2 \rfloor$ this holds exactly when
$m \equiv 0$ or $3 \pmod 4$. The restriction belongs to the \emph{balanced split}, not
to the scheme: choosing $\mathsf{M1}$ freely makes any $m$ admissible --- $m = 10$ at
$\mathsf{M1} = 4$ has a row half of $6$.

\emph{Drawability.} The $Q$ queries are \emph{distinct} indices into the folded domain,
so $Q \le 2^{\mathsf{RowVars} - \ell + \mathsf{LogRate}}$. At $Q = 43$ this is what
rules out the small sizes that pass the parity rule.

\section{Deliberate departures from the paper}
\label{sec:notdone}

The implementation is a faithful but partial instantiation. What is missing, and why it
matters:

\begin{description}[leftmargin=*,style=nextline]
  \item[Not WHIR proper.] The reduced polynomial is sent in plain rather than recursed
    on, so the verifier is $O(2^{\mathsf{RowVars}-\ell})$ rather than
    polylogarithmic. Soundness is unaffected --- sending more is not sending less ---
    but proof size and verifier time are.
  \item[No $O(n^{1/4})$ layer.] The paper's second folding layer runs over the
    \emph{generator} oracle \paper{\S\,group-optimizations}. It is absent, which is
    precisely why $\mathsf{M1}$ defaults to $\lfloor m/2 \rfloor$ here where the
    reference implementation uses $\lfloor m/2 \rfloor - 2$: without leg-2 folding,
    every extra column variable is paid in full.
  \item[No CSP folding.] Related to the above: leg 2 is linear. This is also why the
    reference's larger query count does not apply here --- it accounts for Johnson's
    radius \emph{and} the soundness error of folding CSP, neither of which is in play.
  \item[No zero knowledge.] The proof reveals $Q$ cosets of the oracle and the reduced
    polynomial in plain. This is a commitment scheme, not a ZK argument.
  \item[No serialization, no batching, no trusted setup.] Generators are supplied by
    the caller; proofs are in-memory structures; one evaluation point per proof.
\end{description}

\section{References}

Section names above refer to the Titan paper (eprint version): \S\,titan-overview,
\S\,titan-setup, \S\,titan-commit, \S\,titan-eval-proof, \S\,group-sum-check,
\S\,whir-oracle-algo, \S\,csp-speedup, \S\,partial-eval-csp, \S\,check-C,
\S\,group-optimizations, \S\,titan-pcs-costs, Thm.\,titan-pcs, Lem.\,csp-soundness,
Lem.\,eqind.

The interface this construction is exposed through is documented in
\texttt{docs/crypto/titan-pcs-api.md}; the implementation design record, including the
known defects that no functional test detects, is in \texttt{docs/crypto/titan.md}.

\end{document}
